\scanpage{192}
cas seulement, on peut choisir comme générateur de $B$ sur $A$
\[
\Delta=\tau-\tau'
\]
soumis à la relation
\[
\Delta^2=\Sigma^2-4\Pi\in A
\]
On aura
\[
\tau=\frac12(\Delta+\Sigma)
\]
\[
\tau'=\frac12(-\Delta+\Sigma)
\]
Mais en car. $2$, on aura
\[
\Delta=\tau+\tau'=\Sigma\in A\qquad(\text{en car. }2)
\]
donc $\Delta$ ne risque pas d'engendrer $B$ sur $A$ en car. $2$, donc pas non plus si $2$ non inv. dans $k$ i.e.
\[
k/2k\ne0.
\]
\figure{\begin{tikzpicture}\draw (0,0)--(2.1,0);\end{tikzpicture}}

Cas de quatre variables $x,y,z,t$.
Soit
\[
B_4(\mathcal{A})\subset\operatorname{Sym}(\check{\mathcal{A}}^{4})
\]
l'anneau des fonctions \unsure{polynômes} $F(x,y,z,t)$
\margin{191}
